\section{\texorpdfstring{从方程到矩阵：\(Ax=b\) 的三种读法}{从方程到矩阵：Ax=b 的三种读法}}
\label{sec:03-Linear-Algebra/01-Linear-Systems-and-Matrices/01}

\subsection{一个食堂里的具体困境}

假设你在经营一个小食堂。今天的菜单只有两种套餐：

\begin{itemize}
\tightlist
\item 套餐甲：1 份米饭 + 2 杯饮料，售价 5 元；
\item 套餐乙：3 份米饭 - 1 杯饮料（老板送了你一杯饮料券），售价 4 元。
\end{itemize}

你想知道米饭和饮料各自的单价。设米饭单价 \(x\) 元，饮料单价 \(y\) 元：

\[
\begin{aligned}
x+2y&=5,\\
3x-y&=4.
\end{aligned}
\]

中学消元：第一式给出 \(x=5-2y\)，代入第二式得 \(3(5-2y)-y=4\)，即 \(15-7y=4\)，得 \(y=\frac{11}{7}\)，回代得 \(x=\frac{13}{7}\)。两行就算出来了，没什么压力。

现在食堂扩张，加入了汤：

\begin{itemize}
\tightlist
\item 套餐甲：1 米饭 + 2 饮料 + 1 汤 = 8 元；
\item 套餐乙：3 米饭 - 1 饮料 + 2 汤 = 11 元；
\item 套餐丙：2 米饭 + 1 饮料 + 3 汤 = 13 元。
\end{itemize}

你来试试消元。从第一个方程解出 \(x=8-2y-z\)，代入后两个……等一下，这才三个变量就已经有点乱了。你可能会想：如果食堂有 50 种食材和 50 种套餐，中学那套逐一代入的方法还能用吗？

能，但你会迷失在符号的海洋里。问题出在：变量名 \(x,y,z\) 根本只是占位符，真正干活的是那些系数和右端的数字。每次消元，你做的其实是把系数拿过来算来算去，变量名只是跟着走个过场。

\subsection{把变量名扔掉：矩阵记号从哪里来}

仔细观察上面的三元系统。如果你把系数单独抽出来，排成一张表：

\[
\begin{pmatrix}
1 & 2 & 1 \\
3 & -1 & 2 \\
2 & 1 & 3
\end{pmatrix},
\]

把未知量排成一列：

\[
\begin{pmatrix}x\\ y\\ z\end{pmatrix},
\]

把右端也排成一列：

\[
\begin{pmatrix}8\\ 11\\ 13\end{pmatrix},
\]

那么整个系统可以紧凑地写成

\[
\underbrace{\begin{pmatrix}1&2&1\\3&-1&2\\2&1&3\end{pmatrix}}_{A}
\underbrace{\begin{pmatrix}x\\ y\\ z\end{pmatrix}}_{x}
=
\underbrace{\begin{pmatrix}8\\ 11\\ 13\end{pmatrix}}_{b}.
\]

再把系数矩阵简记为 \(A\)，未知向量简记为 \(x\)，右端向量简记为 \(b\)，就有了那个在整个线性代数中无处不在的等式：

\[
Ax=b.
\]

这个写法最初看起来只是一个缩略符号，但它真正的力量在于：同一个等式 \(Ax=b\) 同时支持三种完全不同的理解方式。这三种读法不是同一件事的不同说法——它们对应不同的问题，揭示不同的结构，各自在某些场景中更有用。

\subsection{行视角：每条规则画出一个几何对象}

回到最简单的二元情况：

\[
\begin{aligned}
x+2y&=5,\\
3x-y&=4.
\end{aligned}
\]

把 \(A\) 的每一行看作一个约束。第一行 \(x+2y=5\) 在平面上是一条直线，第二行 \(3x-y=4\) 是另一条直线。两条直线在平面上相交，那个交点 \((13/7, 11/7)\) 就是解。

推广到三元情况：每个方程是三维空间中的一个平面。三个平面可能刚好交于一点（唯一解），可能沿一条直线相交（无穷多解），也可能三者围出一个空洞（无解——试想三面墙角向外敞开的情形）。

再往高维走，每个方程就是 \(n\) 维空间中的一个超平面。行视角把解方程组理解成寻找几何对象的公共交点。

但你很快会碰到一个麻烦：只有两三个方程时画图还行，十个方程的时候你画不出来。而且，仅凭数方程个数和未知量个数，你无法判断属于哪种情况。两个方程可能说的是同一件事（比如第二个方程只是第一个的两倍），十个方程也可能彼此完全一致。真正重要的是其中有多少\emph{独立的}信息。

这就是行视角的边界：它几何直观，但在高维中无法直接操作。它告诉你解是交点，但没告诉你这个交点怎么算出来，更没告诉你什么样的右端 \(b\) 才可能有解。

\subsection{列视角：用给定的方向合成目标}

看看我们用过的那个二元系统，但换一种打开方式。不把每一行看作一个方程，而是把每一列看作一个方向向量：

\[
x\begin{pmatrix}1\\3\end{pmatrix}
+y\begin{pmatrix}2\\-1\end{pmatrix}
=\begin{pmatrix}5\\4\end{pmatrix}.
\]

这不再是求两条直线的交点。问题变成了：有没有一组系数 \((x,y)\)，使得两个列向量的线性组合刚好等于目标向量 \((5,4)^{\trans}\)？

你可能会想：这不就是换了个说法吗？代数上确实等价，但问题变了——从``找交点''变成了``找组合系数''。

试一下：这两列能组合出 \((5,4)^{\trans}\) 吗？解出来 \(x=13/7,y=11/7\)，也就是说

\[
\frac{13}{7}\begin{pmatrix}1\\3\end{pmatrix}
+\frac{11}{7}\begin{pmatrix}2\\-1\end{pmatrix}
=\begin{pmatrix}5\\4\end{pmatrix}.
\]

两列不是彼此的倍数（一列不落在另一列的方向上），它们可以张成整个二维平面，所以任意二维目标都能被精确合成。

来看一个不同的例子：

\[
\begin{pmatrix}1&2\\2&4\end{pmatrix}
\begin{pmatrix}x\\y\end{pmatrix}
=\begin{pmatrix}3\\6\end{pmatrix}.
\]

两列分别是 \((1,2)^{\trans}\) 和 \((2,4)^{\trans}\)——第二列刚好是第一列的两倍。这两列都落在同一条直线上，它们只能合成这条直线上的向量。而右端 \((3,6)^{\trans}\) 恰好也在这条直线上（\(=3\times(1,2)^{\trans}\)），所以有解。但换成 \((3,7)^{\trans}\) 就没辙了——它不在那条直线上。

一般地，矩阵向量乘法的列形式是

\[
Ax
=x_1a_1+\cdots+x_na_n,
\]

其中 \(a_j\) 是 \(A\) 的第 \(j\) 列。所有可能的 \(Ax\) 构成的集合称为 \(A\) 的\emph{列空间}，记作 \(\operatorname{Col}(A)\)。于是

\[
Ax=b\text{ 有解}
\quad\Longleftrightarrow\quad
b\in\operatorname{Col}(A).
\]

这是贯穿整门线性代数的核心观察。你不需要解方程就能判断可解性——只要看 \(b\) 是否落在列空间里。后面遇到无解系统时，我们也正是把 \(b\) 投影到列空间中最近的点上去寻找最佳近似解。

列视角的威力在于：它把``有没有解''变成了一个关于``方向够不够''的几何判断，与方程的个数、变量的个数没有直接关系。

\subsection{变换视角：矩阵是一个函数}

再换一个角度看你写的 \(Ax=b\)。不要把它看作方程组，也不看作列组合——把它看作：左边输入了一个 \(x\)，右边输出了一个 \(b\)。

矩阵 \(A\in\R^{m\times n}\) 定义了一个映射

\[
T_A:\R^n\to\R^m,
\qquad
T_A(x)=Ax.
\]

它满足

\[
A(\alpha u+\beta v)=\alpha Au+\beta Av,
\]

所以是\emph{线性变换}：输入缩放了，输出就缩放；两个输入相加，输出也相加。这个性质极其强大——它意味着你只需要知道这个变换在几个基本方向上做了什么，就能知道它在任意方向上的效果。

\(n\) 维输入空间有 \(n\) 个标准基方向。第 \(j\) 个标准基向量 \(e_j\)（第 \(j\) 位为 1、其余为 0 的向量）经过 \(A\) 后变成什么？直接乘一下：

\[
Ae_j=a_j.
\]

\(A\) 的第 \(j\) 列！任意输入 \(x=\sum_j x_je_j\) 经过变换：

\[
Ax=A\Bigl(\sum_j x_je_j\Bigr)=\sum_j x_j(Ae_j)=\sum_j x_ja_j.
\]

矩阵的每一列记录了这个线性变换把一个基向量送到了哪里。正因为变换是线性的，知道各列就完全确定了整个变换。

例如

\[
A=
\begin{pmatrix}
1&2\\
0&-1\\
3&1
\end{pmatrix}
\]

把二维输入 \((x,y)^{\trans}\) 送到三维输出：

\[
Ax=
\begin{pmatrix}
x+2y\\-y\\3x+y
\end{pmatrix}.
\]

第一列告诉你：输入方向的 \((1,0)^{\trans}\) 被送到了输出空间的 \((1,0,3)^{\trans}\)。第二列告诉你：\((0,1)^{\trans}\) 被送到了 \((2,-1,1)^{\trans}\)。尺寸 \(3\times2\) 的含义：2 列对应 2 维输入，3 行对应 3 维输出。

变换视角统一了列视角：列空间就是变换的值域（所有可能的输出）；它也自然地引出了下一个问题——两个变换先后做，效果怎么算？

\subsection{矩阵乘法的真正动机：变换的复合}

假设你有一个二维输入，先经过 \(B:\R^2\to\R^3\)，再经过 \(A:\R^3\to\R^2\)：

\[
x\longmapsto Bx\longmapsto A(Bx).
\]

整体效果还是一个从 \(\R^2\) 到 \(\R^2\) 的映射。两段都是线性的，整体自然也是线性的。那么，整体变换对应的矩阵是什么？

取具体的矩阵来试：

\[
B=
\begin{pmatrix}
1&0\\
0&1\\
1&1
\end{pmatrix},\qquad
A=
\begin{pmatrix}
1&2&0\\
0&-1&3
\end{pmatrix}.
\]

输入 \((x,y)^{\trans}\)，先得

\[
Bx=
\begin{pmatrix}
x\\y\\x+y
\end{pmatrix},
\]

再经 \(A\)：

\[
A(Bx)=
\begin{pmatrix}
1\cdot x+2\cdot y+0\cdot (x+y)\\
0\cdot x+(-1)\cdot y+3\cdot (x+y)
\end{pmatrix}
=
\begin{pmatrix}
x+2y\\
3x+2y
\end{pmatrix}.
\]

整体效果相当于一个矩阵 \(C\) 作用于 \((x,y)^{\trans}\)：

\[
C=
\begin{pmatrix}
1&2\\
3&2
\end{pmatrix}.
\]

\(C\) 的列是什么？第一列是 \(A(B\,e_1)=A\,b_1\)，第二列是 \(A(B\,e_2)=A\,b_2\)。其中 \(b_j\) 是 \(B\) 的第 \(j\) 列。计算：

\[
Ab_1=
\begin{pmatrix}
1&2&0\\
0&-1&3
\end{pmatrix}
\begin{pmatrix}1\\0\\1\end{pmatrix}
=\begin{pmatrix}1\\3\end{pmatrix},
\qquad
Ab_2=
\begin{pmatrix}
1&2&0\\
0&-1&3
\end{pmatrix}
\begin{pmatrix}0\\1\\1\end{pmatrix}
=\begin{pmatrix}2\\2\end{pmatrix}.
\]

正是 \(C\) 的两列。一般地：

\[
AB=\begin{pmatrix}Ab_1&\cdots&Ab_p\end{pmatrix}.
\]

逐元素写出第 \(i\) 行第 \(j\) 列：

\[
(AB)_{ij}=\sum_{k=1}^n a_{ik}b_{kj}.
\]

这就是``行乘列''公式的来历。它不是什么任意的定义——它是``先做 \(B\) 后做 \(A\)''这个自然操作在坐标下的表达。求和指标 \(k\) 对应中间空间 \(\R^n\)，所以只有 \(A\) 的列数等于 \(B\) 的行数时乘法才有定义：

\[
(m\times n)(n\times p)=(m\times p).
\]

如果你以前只是背下了行乘列的规则，现在可以松一口气：它不过是变换复合之后，逐列计算新矩阵的列而已。记住``\(AB\) 的第 \(j\) 列 = \(A\) 乘以 \(B\) 的第 \(j\) 列''，比死记求和公式管用得多。

\subsection{顺序不能随便换：一个几何实验}

\(AB\) 表示先做 \(B\)，后做 \(A\)。如果交换顺序变成 \(BA\)，就算尺寸碰巧对得上，结果通常也不同。

动手试一个例子。取平面上一个点 \((1,1)^{\trans}\)，先对它做水平方向的剪切（\(x\) 坐标加上 \(y\) 坐标），再做水平拉伸两倍。

剪切矩阵：
\[
B=\begin{pmatrix}1&1\\0&1\end{pmatrix},
\]
拉伸矩阵：
\[
A=\begin{pmatrix}2&0\\0&1\end{pmatrix}.
\]

先剪切再拉伸：

\[
AB=\begin{pmatrix}2&0\\0&1\end{pmatrix}
\begin{pmatrix}1&1\\0&1\end{pmatrix}
=\begin{pmatrix}2&2\\0&1\end{pmatrix}.
\]

作用于 \((1,1)^{\trans}\) 得 \((4,1)^{\trans}\)。

反过来，先拉伸再剪切：

\[
BA=\begin{pmatrix}1&1\\0&1\end{pmatrix}
\begin{pmatrix}2&0\\0&1\end{pmatrix}
=\begin{pmatrix}2&1\\0&1\end{pmatrix}.
\]

作用于 \((1,1)^{\trans}\) 得 \((3,1)^{\trans}\)。

结果不同，因为先剪切时把 \(y\) 坐标混进了 \(x\) 坐标，后面的拉伸把这个混合后的 \(x\) 放大了两倍；先拉伸时剪切还没发生，拉伸只放大原始的 \(x\)，随后剪切才加入 \(y\) 的贡献。操作的顺序改变了坐标被修改的时机。

你可能会问：有没有什么时候 \(AB=BA\)？有，比如两个对角矩阵（各自独立缩放各维），或者 \(A\) 是单位矩阵，或者 \(A\) 和 \(B\) 都是同一个矩阵的幂。但一般而言，不可交换是常态。

矩阵乘法满足结合律 \((AB)C=A(BC)\)（两边都表示按相同顺序复合三个变换）和分配律，但不满足交换律。还不能从 \(AB=AC\) 随意把左边的 \(A\)``约掉''——什么时候允许这样做，要等到逆矩阵部分才能说清楚。

\subsection{转置：把行列翻过来在干什么}

矩阵转置 \(A^{\trans}\) 把行变成列、列变成行。若 \(A\) 是 \(m\times n\)，则 \(A^{\trans}\) 是 \(n\times m\)。

这个操作看起来只是把数字重新排列了一下，但它有一条不那么显然的性质：

\[
(AB)^{\trans}=B^{\trans}A^{\trans}.
\]

为什么会反转顺序？回想 \(AB\) 的意义：先做 \(B\) 后做 \(A\)。转置之后，输入和输出的配对关系反向传递——原来从输入到 \(B\) 再到 \(A\) 到输出，转置后从``新输入''（原输出空间）到 \(A^{\trans}\) 再到 \(B^{\trans}\) 到``新输出''（原输入空间）。次序反转是方向反转的自然结果。

用具体的数验证一下：

\[
A=\begin{pmatrix}1&2\\3&4\end{pmatrix},\;
B=\begin{pmatrix}5&6\\7&8\end{pmatrix},\;
AB=\begin{pmatrix}19&22\\43&50\end{pmatrix},\;
(AB)^{\trans}=\begin{pmatrix}19&43\\22&50\end{pmatrix}.
\]

\[
B^{\trans}A^{\trans}=
\begin{pmatrix}5&7\\6&8\end{pmatrix}
\begin{pmatrix}1&3\\2&4\end{pmatrix}
=\begin{pmatrix}19&43\\22&50\end{pmatrix}.
\]

逐元素也可以一般地验证：

\[
\bigl((AB)^{\trans}\bigr)_{ij}=(AB)_{ji}
=\sum_{k}a_{jk}b_{ki}
=\sum_{k}(B^{\trans})_{ik}(A^{\trans})_{kj}
=(B^{\trans}A^{\trans})_{ij}.
\]

转置在后面讨论正交矩阵、最小二乘和对称矩阵时会反复出现。现在你只需要记住它的几何直觉：转置把变换的``输入输出方向''反过来。

\subsection{三种读法，一个等式}

回到 \(Ax=b\)。你现在手里有三副眼镜：

\begin{itemize}
\tightlist
\item 行视角：\(b\) 的每个分量来自一个方程（约束），解是这些约束几何对象的交点；
\item 列视角：\(b\) 是 \(A\) 各列的线性组合，组合系数就是解 \(x\)；
\item 变换视角：\(A\) 把输入 \(x\) 映射到输出 \(b\)，解是使映射到达 \(b\) 的那个输入。
\end{itemize}

三种读法不需要你选一个——它们同时成立，互相补充。行视角最直观，列视角揭示了可解性的核心判据，变换视角让你把矩阵当作可以复合、可以求逆的操作对象。接下来讨论消元、逆矩阵和向量空间时，我们会根据具体问题在三种视角之间来回切换。

延伸阅读：MIT 18.06：The Geometry of Linear Equations。
